\documentclass[10pt,letterpaper]{article}
\usepackage{amsmath,amssymb,amsthm}
\setlength{\parindent}{0pt}
\setlength{\parskip}{6pt}
\begin{document}
\section{Linear algebra review}
Let $A$ be an $n\times n$ real matrix. A scalar $\lambda$ is an \emph{eigenvalue} if $Av=\lambda v$ for some nonzero vector $v$.

\begin{itemize}
\item The eigenvalues are the roots of $\det(A-\lambda I)=0$.
\item The trace of $A$ is the sum of the eigenvalues, counted with multiplicity.
\item The determinant is their product.
  \begin{itemize}
  \item In particular $A$ is invertible exactly when $0$ is not an eigenvalue.
  \item For symmetric $A$ all eigenvalues are real.
    \begin{itemize}
    \item Eigenvectors for distinct eigenvalues are orthogonal.
      \begin{itemize}
      \item This gives the spectral theorem.
      \end{itemize}
    \end{itemize}
  \end{itemize}
\end{itemize}

\subsection{A $2\times2$ example}
For $A=\begin{pmatrix}2&1\\1&2\end{pmatrix}$ we find
\begin{gather*}
\det(A-\lambda I)=(2-\lambda)^2-1=\lambda^2-4\lambda+3,\\
\lambda_1=1,\quad\lambda_2=3.
\end{gather*}
\end{document}
